\begin{abstract}
%Pit thermal energy storage (PTES) can provide district-heating flexibility by shifting heat supply across time, but existing tariffs may provide insufficient incentives for its utilisation. Under a monthly tariff, local operators do not observe hourly variations in upstream production costs. Charging storage therefore offers little advantage over direct heat imports, while thermal losses further discourage cycling.

Pit thermal energy storage (PTES) can provide district-heating flexibility by shifting heat production relative to heat demand. In systems with combined heat and power (CHP) generation, it allows heat to be produced and stored when electricity prices are high and the value of co-produced electricity makes CHP operation particularly attractive. However, existing tariff structures may not transmit this time-varying system value to local storage operators. Under a monthly tariff, hourly variations in upstream heat-production costs are not reflected in the price paid by the local operator. Consequently, the tariff provides little economic incentive to charge and discharge the storage in response to these variations, while thermal losses further discourage cycling.

This paper develops a method for determining a dynamic heat tariff and evaluates whether it improves PTES utilisation. The hourly tariff is derived from the marginal heat cost of an upstream biomass combined heat and power plant, accounting for fuel and operating costs, carbon costs, conversion efficiencies, and the market value of co-produced electricity. % A fixed/revenue adjustment component supports cost recovery without changing the marginal dispatch signal. (Possible idea, lets see) 

The tariff is soft-linked to an hourly mixed integer model of the Høje Taastrup district heating system, jointly optimising direct heat imports, a 3.3~GWh PTES, and two heat pumps under realistic operating constraints. Using 2024 heat-demand and electricity price data, the dynamic tariff is compared with the incumbent monthly VEKS tariff and an unconstrained PTES benchmark.

%\textbf{Include another paragraph about the results}
\end{abstract}
